\documentclass[12pt,aspectratio=169,ignorenonframetext]{ltxbeamer}

\usepackage[spanish,es-nodecimaldot,es-tabla]{babel}

\title{Introducción a {\lmr\TikZ}}
\subtitle{Curso de {\lmr\LaTeX}}
\author{Joar Buitrago\texorpdfstring{\thanks{jebuitragoc@unal.edu.co}}{ <jebuitragoc@unal.edu.co>}}
\date{}

\addbibresource{bibliography.bib}


\newcommand{\TikZ}{Ti\textit{k}Z}


\begin{document}
\begin{frame}
\titlepage
\end{frame}



\section{¿Qué es {\lmr\TikZ}?}

\begin{frame}
\frametitle{¿Qué es {\lmr\TikZ}?}

\pause

\begin{itemize}[<+->]
    \item {\lmr\TikZ} es una capa de macros de {\lmr PGF}.
    \item Nos permite realizar dibujos directamente en {\lmr\LaTeX}.
    \item {\lmr\TikZ} genera dibujos en vectoriales, que se acomodan perfectamente al texto circundante.
\end{itemize}

\bigskip

\begin{block}<+->{Dato curioso}
    \lmr
    \TikZ{} ist \textit{kein} Zeichenprogramm

    \TikZ{} no es un programa de dibujo
\end{block}
\end{frame}





\section{Fundamentos}

\begin{frame}[fragile]
\begin{center}
    \resizebox{0.8\linewidth}{!}{\mintinline{latex}{\usepackage{tikz}}}
\end{center}
\end{frame}


\begin{frame}[fragile]
\begin{minted}{latex}
\begin{tikzpicture}
    <Dibujo...>
\end{tikzpicture}
\end{minted}
\end{frame}

\begin{frame}[fragile]
\begin{minted}{latex}
\tikz{<Dibujo ...>}
\end{minted}
\end{frame}

\begin{frame}
\frametitle{Sistema de Coordenadas y Unidades}

\begin{itemize}
\item \textbf{Sistema Cartesiano Estándar}: El punto \texttt{(0,0)} es el origen.
\item \textbf{Unidades}: Por defecto, la unidad es \texttt{cm}.
\end{itemize}
\end{frame}

\begin{frame}[fragile]
\frametitle{Sistema de Coordenadas y Unidades}
\framesubtitle{Coordenadas Absolutas}

\begin{columns}
    \begin{column}{0.45\textwidth}
        \begin{block}{}
            \begin{minted}{latex}
\begin{tikzpicture}
    \draw (0,0) -- (4,0);
    \draw (0,0) -- (0,3);
\end{tikzpicture}
            \end{minted}
        \end{block}
    \end{column}

    \begin{column}{0.1\textwidth}
        \centering $\rightarrow$
    \end{column}

    \begin{column}{0.45\textwidth}
        \begin{block}{}
            \begin{tikzpicture}[very thick]
                \draw (0,0) -- (4,0);
                \draw (0,0) -- (0,3);
            \end{tikzpicture}
        \end{block}
    \end{column}
\end{columns}
\end{frame}


\begin{frame}[fragile]
\frametitle{Sistema de Coordenadas y Unidades}
\framesubtitle{Coordenadas Relativas}

Definen puntos en relación con el punto anterior, ideal para secuencias.

\begin{columns}
    \begin{column}{0.45\textwidth}
        \begin{block}{}
            \begin{minted}{latex}
\begin{tikzpicture}
    \draw (0,0) -- ++(4,0)
        -- +(0,2) -- ++(-2,0);
\end{tikzpicture}
            \end{minted}
        \end{block}
    \end{column}

    \begin{column}{0.1\textwidth}
        \centering $\rightarrow$
    \end{column}

    \begin{column}{0.45\textwidth}
        \begin{block}{}
            \begin{tikzpicture}[very thick]
                \draw (0,0) -- ++(4,0) -- +(0,2) -- ++(-2,0);
            \end{tikzpicture}
        \end{block}
    \end{column}
\end{columns}
\pause

\begin{description}
    \item[+] Ir y volver
    \item[++] Ir y quedarse
\end{description}
\end{frame}



\begin{frame}[fragile]
\frametitle{Sistema de Coordenadas y Unidades}
\framesubtitle{cycle}

La palabra clave \texttt{cycle} cierra la ruta hasta el punto inicial.

\begin{columns}
    \begin{column}{0.45\textwidth}
        \begin{block}{}
            \begin{minted}{latex}
\begin{tikzpicture}
    \draw (0,0) -- ++(4,0)
        -- +(0,2) -- cycle;
\end{tikzpicture}
            \end{minted}
        \end{block}
    \end{column}

    \begin{column}{0.1\textwidth}
        \centering $\rightarrow$
    \end{column}

    \begin{column}{0.45\textwidth}
        \begin{block}{}
            \begin{tikzpicture}[very thick]
                \draw (0,0) -- ++(4,0) -- +(0,2) -- cycle;
            \end{tikzpicture}
        \end{block}
    \end{column}
\end{columns}
\end{frame}





\begin{frame}[fragile]
\frametitle{Sistema de Coordenadas y Unidades}
\framesubtitle{Unidades}

Podemos usar cualquier dimensión de {\lmr\TeX} para dibujar.

\begin{center}
\scriptsize
\begin{tabular}{>{\ttfamily}lll}
    \toprule
    Símbolo & Unidad & Definición \\
    \midrule
    in & Pulgada & \\
    cm & Centímetro & 2.54 cm = 1 in \\
    mm & Milimetro & 10 mm = 1 cm \\
    pt & Punto & 72.27 pt = 1 in \\
    pc & Pica & 1 pc = 12 pt \\
    bp & Punto Grande & 72 bp = 1 in \\
    dd & Punto Didot & 1157 dd = 1238 pt \\
    cc & Cicero & 1 cc = 12 dd \\
    sp & Punto Escalado & 65536 sp = 1 pt \\
    ex & Equis & El alto de la \texttt{x} en la fuente actual \\
    em & Eme & El ancho de la \texttt{M} en la fuente actual \\
    fil & Infinito de Primer Orden & \(\aleph_0\) \\
    fill & Infinito de Segundo Orden & \(\aleph_1\) \\
    filll & Infinito de Tercer Orden & \(\aleph_2\) \\
    \bottomrule
\end{tabular}
\end{center}
\end{frame}

\begin{frame}[fragile]
\frametitle{Sistema de Coordenadas y Unidades}
\framesubtitle{Unidades}

\begin{columns}
    \begin{column}{0.45\textwidth}
        \begin{block}{}
            \begin{minted}{latex}
\begin{tikzpicture}
    \draw (0,0) -- (0,4ex);
\end{tikzpicture}
            \end{minted}
        \end{block}
    \end{column}

    \begin{column}{0.1\textwidth}
        \centering $\rightarrow$
    \end{column}

    \begin{column}{0.45\textwidth}
        \begin{block}{}
            \begin{tikzpicture}[very thick]
                \draw (0,0) -- (0,4ex);
                \node at (1.0ex,0ex) [inner sep=0,anchor=base] {x};
                \node at (1.3ex,1ex) [inner sep=0,anchor=base] {x};
                \node at (1.0ex,2ex) [inner sep=0,anchor=base] {x};
                \node at (1.3ex,3ex) [inner sep=0,anchor=base] {x};
            \end{tikzpicture}
        \end{block}
    \end{column}
\end{columns}
\end{frame}


\begin{frame}[fragile]
\frametitle{Sistema de Coordenadas y Unidades}
\framesubtitle{Coordenadas Polares}

Para {\lmr\TikZ}, el primer argumento es el \alert{ángulo} y el segundo el \alert{radio}.
Por defecto usa grados, pero con el sufijo \texttt{r} usamos radianes.

\pause

\begin{columns}
    \begin{column}{0.45\textwidth}
        \begin{block}{}
            \begin{minted}{latex}
\begin{tikzpicture}
    \draw (0,0) -- (15:20ex);
    \draw (0,0) -- (1r:4em);
    \draw (0,0) -- (pi r:1in);
\end{tikzpicture}
            \end{minted}
        \end{block}
    \end{column}

    \begin{column}{0.1\textwidth}
        \centering $\rightarrow$
    \end{column}

    \begin{column}{0.45\textwidth}
        \begin{block}{}
            \begin{tikzpicture}[very thick]
                \draw (0,0) -- (15:20ex);
                \draw (0,0) -- (1r:4em);
                \draw (0,0) -- (pi r:1in);
            \end{tikzpicture}
        \end{block}
    \end{column}
\end{columns}
\end{frame}



\section{Primitivas de Dibujo}

\begin{frame}[fragile]
\frametitle{Primitivas de Dibujo}
\framesubtitle{Líneas}
El operador \texttt{--} define una línea recta.

\begin{columns}
    \begin{column}{0.45\textwidth}
        \begin{block}{}
            \begin{minted}{latex}
\begin{tikzpicture}
    \draw (0,0) -- (1,1)
        -- (2,0) -- (3,1);
\end{tikzpicture}
            \end{minted}
        \end{block}
    \end{column}

    \begin{column}{0.1\textwidth}
        \centering $\rightarrow$
    \end{column}

    \begin{column}{0.45\textwidth}
        \begin{block}{}
            \begin{tikzpicture}[very thick]
                \draw (0,0) -- (1,1) -- (2,0) -- (3,1);
            \end{tikzpicture}
        \end{block}
    \end{column}
\end{columns}
\end{frame}


\begin{frame}[fragile]
\frametitle{Primitivas de Dibujo}
\framesubtitle{Curvas de Bézier}
El operador \texttt{..} define una curva de Bézier, junto con sus puntos de control.

\begin{columns}
    \begin{column}{0.45\textwidth}
        \begin{block}{}
            \begin{minted}{latex}
\begin{tikzpicture}
    \draw (0,0)
        .. controls
            (0,2) and (2,2)
        .. (2,0);
\end{tikzpicture}
            \end{minted}
        \end{block}
    \end{column}

    \begin{column}{0.1\textwidth}
        \centering $\rightarrow$
    \end{column}

    \begin{column}{0.45\textwidth}
        \begin{block}{}
            \begin{onlyenv}<1>
                \begin{tikzpicture}[very thick]
                    \draw (0,0) .. controls (0,2) and (2,2) .. (2,0);
                    \draw (0,0) node [inner sep=2pt] {};
                    \draw (2,0) node [inner sep=2pt] {};
                    \draw (0,2) node [inner sep=2pt, circle] {};
                    \draw (2,2) node [inner sep=2pt, circle] {};
                \end{tikzpicture}
            \end{onlyenv}
            \begin{onlyenv}<2->
                \begin{tikzpicture}[very thick]
                    \draw (0,0) .. controls (0,2) and (2,2) .. (2,0);

                    \draw[dashed, help lines, thick] (0,0) node [inner sep=2pt, fill] {} -- (0,2) node [fill, inner sep=2pt, circle] {};
                    \draw[dashed, help lines, thick] (2,0) node [inner sep=2pt, fill] {} -- (2,2) node [fill, inner sep=2pt, circle] {};
                \end{tikzpicture}
            \end{onlyenv}
        \end{block}
    \end{column}
\end{columns}
\end{frame}



\begin{frame}[fragile]
\frametitle{Primitivas de Dibujo}
\framesubtitle{Rectángulo}
El operador \texttt{rectangle} define un rectángulo.

\begin{columns}
    \begin{column}{0.45\textwidth}
        \begin{block}{}
            \begin{minted}{latex}
\begin{tikzpicture}
    \draw (0,0) rectangle (3,1);
\end{tikzpicture}
            \end{minted}
        \end{block}
    \end{column}

    \begin{column}{0.1\textwidth}
        \centering $\rightarrow$
    \end{column}

    \begin{column}{0.45\textwidth}
        \begin{block}{}
            \begin{tikzpicture}[very thick]
                \draw (0,0) rectangle (3,1);
            \end{tikzpicture}
        \end{block}
    \end{column}
\end{columns}
\end{frame}



\begin{frame}[fragile]
\frametitle{Primitivas de Dibujo}
\framesubtitle{Círculo}
El operador \texttt{circle} define un círculo, cuyo argumento es el radio.

\begin{columns}
    \begin{column}{0.45\textwidth}
        \begin{block}{}
            \begin{minted}{latex}
\begin{tikzpicture}
    \draw (0,0) circle (2);
\end{tikzpicture}
            \end{minted}
        \end{block}
    \end{column}

    \begin{column}{0.1\textwidth}
        \centering $\rightarrow$
    \end{column}

    \begin{column}{0.45\textwidth}
        \begin{block}{}
            \begin{tikzpicture}[very thick]
                \draw (0,0) circle (2);
            \end{tikzpicture}
        \end{block}
    \end{column}
\end{columns}
\end{frame}



\begin{frame}[fragile]
\frametitle{Primitivas de Dibujo}
\framesubtitle{Elipse}
El operador \texttt{ellipse} define una elipse, cuyo argumentos son los radios.

\begin{columns}
    \begin{column}{0.45\textwidth}
        \begin{block}{}
            \begin{minted}{latex}
\begin{tikzpicture}
    \draw (0,0)
        ellipse (2 and 1);
\end{tikzpicture}
            \end{minted}
        \end{block}
    \end{column}

    \begin{column}{0.1\textwidth}
        \centering $\rightarrow$
    \end{column}

    \begin{column}{0.45\textwidth}
        \begin{block}{}
            \begin{tikzpicture}[very thick]
                \draw (0,0) ellipse (2 and 1);
            \end{tikzpicture}
        \end{block}
    \end{column}
\end{columns}
\end{frame}



\begin{frame}[fragile]
\frametitle{Primitivas de Dibujo}
\framesubtitle{Arco}
El operador \texttt{arc} define un arco, cuyo argumentos son el ángulo inicial, el ángulo final y los radios.

\begin{columns}
    \begin{column}{0.45\textwidth}
        \begin{block}{}
            \begin{minted}{latex}
\begin{tikzpicture}
    \draw (0,0)
        arc (10:130:2 and 1);
\end{tikzpicture}
            \end{minted}
        \end{block}
    \end{column}

    \begin{column}{0.1\textwidth}
        \centering $\rightarrow$
    \end{column}

    \begin{column}{0.45\textwidth}
        \begin{block}{}
            \begin{onlyenv}<1>
                \begin{tikzpicture}[very thick]
                    \useasboundingbox (0,0) circle (2 and 1);
                    \draw (10:2 and 1) arc (10:130:2 and 1);
                \end{tikzpicture}
            \end{onlyenv}
            \begin{onlyenv}<2->
                \begin{tikzpicture}[very thick]
                    \useasboundingbox (0,0) circle (2 and 1);
                    \draw[dashed,help lines,thick] (0,0) circle (2 and 1);
                    \draw (10:2 and 1) arc (10:130:2 and 1);
                \end{tikzpicture}
            \end{onlyenv}
        \end{block}
    \end{column}
\end{columns}
\end{frame}



\section{Estilo y Color}

\begin{frame}[fragile]
\frametitle{Opciones de Estilo y Color}

Las opciones se pasan entre corchetes \mintinline{latex}{\draw [opciones]}

\begin{itemize}
    \item \textbf{Grosor}:
    \begin{itemize}
        \item \texttt{thick}
        \item \texttt{ultra thick}
        \item \texttt{thin}
    \end{itemize}
    \item \textbf{Punteado}:
    \begin{itemize}
        \item \texttt{dashed}
        \item \texttt{dotted}
        \item \texttt{loosely dashed}
    \end{itemize}
    \item \textbf{Color}:
    \begin{itemize}
        \item \texttt{red} \only<2->{--- Rojo -- Blanco \tikz{\node [fill,red] {};}}
        \item \texttt{blue!50} \only<3->{--- Azul -- 50\% -- Blanco \tikz{\node [fill,blue!50] {};}}
        \item \texttt{blue!20!green} \only<4->{--- Azul -- 20\% -- Verde -- Blanco \tikz{\node [fill,blue!20!green] {};}}
        \item \texttt{blue!20!green!50} \only<5->{--- Azul -- 20\% -- Verde -- 50\% -- Blanco \tikz{\node [fill,blue!20!green!50] {};}}
    \end{itemize}
\end{itemize}
\end{frame}



\section{Nodos y Conexiones}

\begin{frame}[fragile]
\frametitle{Nodos}

Los Nodos se usan para posicionar texto o bloques.

\begin{columns}
    \begin{column}{0.45\textwidth}
        \begin{block}{}
            \begin{minted}{latex}
\begin{tikzpicture}
    \node {Texto};
\end{tikzpicture}
            \end{minted}
        \end{block}
    \end{column}

    \begin{column}{0.1\textwidth}
        \centering $\rightarrow$
    \end{column}

    \begin{column}{0.45\textwidth}
        \begin{block}{}
            \begin{tikzpicture}
                \node {Texto};
            \end{tikzpicture}
        \end{block}
    \end{column}
\end{columns}
\end{frame}



\begin{frame}[fragile]
\frametitle{Nodos}

Podemos ubicarlos, y desfasarlos.

\begin{columns}
    \begin{column}{0.45\textwidth}
        \begin{block}{}
            \begin{minted}{latex}
\begin{tikzpicture}
    \node at (1,2)
        [below] {Debajo};
    \draw (1,2) node
        [above] {Encima};
    \draw (1,2) node
        [right] {Derecha};
\end{tikzpicture}
            \end{minted}
        \end{block}
    \end{column}

    \begin{column}{0.1\textwidth}
        \centering $\rightarrow$
    \end{column}

    \begin{column}{0.45\textwidth}
        \begin{block}{}
            \begin{tikzpicture}
                \node at (1,2) [below] {Debajo};
                \draw (1,2) node [above] {Encima};
                \draw (1,2) node [right] {Derecha};
            \end{tikzpicture}
        \end{block}
    \end{column}
\end{columns}
\end{frame}


\begin{frame}[fragile]
\frametitle{Nodos}

Podemos darles un borde o rellenarlos.

\begin{columns}
    \begin{column}{0.45\textwidth}
        \begin{block}{}
            \begin{minted}{latex}
\begin{tikzpicture}
    \node at (1,4)
        [draw] {Texto};
    \draw (1,2) node
        [fill] {Texto};
\end{tikzpicture}
            \end{minted}
        \end{block}
    \end{column}

    \begin{column}{0.1\textwidth}
        \centering $\rightarrow$
    \end{column}

    \begin{column}{0.45\textwidth}
        \begin{block}{}
            \begin{tikzpicture}
                \node at (1,4) [draw] {Texto};
                \draw (1,2) node [fill, white] {\color{black} Texto};
            \end{tikzpicture}
        \end{block}
    \end{column}
\end{columns}
\end{frame}


\begin{frame}[fragile]
\frametitle{Nodos}

Podemos nombrar nodos, y referenciarlos.

\begin{columns}
    \begin{column}{0.45\textwidth}
        \begin{block}{}
            \begin{minted}{latex}
\begin{tikzpicture}
    \node (first) {1};
    \node (second)
        [below of=first] {2};
    \node (third)
        [right of=second] {3};
\end{tikzpicture}
            \end{minted}
        \end{block}
    \end{column}

    \begin{column}{0.1\textwidth}
        \centering $\rightarrow$
    \end{column}

    \begin{column}{0.45\textwidth}
        \begin{block}{}
            \begin{tikzpicture}
                \node (first) {1};
                \node (second) [below of=first] {2};
                \node (third) [right of=second] {3};
            \end{tikzpicture}
        \end{block}
    \end{column}
\end{columns}
\end{frame}




\begin{frame}[fragile]
\frametitle{Nodos}

También podemos unirlos, por medio de líneas, u otras primitivas de {\lmr\TikZ}.

\begin{columns}
    \begin{column}{0.45\textwidth}
        \begin{block}{}
            \begin{minted}{latex}
\begin{tikzpicture}
    \node (first) {1};
    \node (second)
        [below of=first] {2};
    \node (third)
        [right of=second] {3};

    \draw[->] (first.south)
        -- (second.north);
    \draw[->] (second.east)
        -- (third.west);
\end{tikzpicture}
            \end{minted}
        \end{block}
    \end{column}

    \begin{column}{0.1\textwidth}
        \centering $\rightarrow$
    \end{column}

    \begin{column}{0.45\textwidth}
        \begin{block}{}
            \begin{tikzpicture}
                \node (first) {1};
                \node (second) [below of=first] {2};
                \node (third) [right of=second] {3};

                \draw[->] (first.south) -- (second.north);
                \draw[->] (second.east) -- (third.west);
            \end{tikzpicture}
        \end{block}
    \end{column}
\end{columns}
\end{frame}




\section{Diagramación}

\begin{frame}[fragile]
\frametitle{Diagramas de Flujo}
\begin{onlyenv}<1>
    \begin{minted}{latex}
\begin{tikzpicture}[draw,node distance=1cm]
    \node (start) {Inicio};
    \node (get) [below of=start] {Obtener};
    \node (check) [below of=get] {¿Válido?};
    \node (process)
        [below right of=check,node distance=2cm]
        {Procesar};

    \draw [->] (start) -- (get);
    \draw [->] (get) -- (check);
    \draw [->] (check) -| (process)
        node [midway,above right] {Sí};
\end{tikzpicture}
    \end{minted}
\end{onlyenv}
\begin{onlyenv}<2->
    \begin{center}
        \begin{tikzpicture}[draw,node distance=1cm]
            \node (start) {Inicio};
            \node (get) [below of=start] {Obtener};
            \node (check) [below of=get] {¿Válido?};
            \node (process) [below right of=check,node distance=2cm] {Procesar};

            \draw [->] (start) -- (get);
            \draw [->] (get) -- (check);
            \draw [->] (check) -| (process) node [midway,above right] {Sí};
        \end{tikzpicture}
    \end{center}
\end{onlyenv}
\end{frame}

\begin{frame}[fragile]
\frametitle{Óptica}
\begin{onlyenv}<1>
    \footnotesize
    \begin{minted}{latex}
\begin{tikzpicture}
    \draw [cyan,<->,very thick] (1,-2) -- (1,2);
    \draw [cyan,>-<,very thick] (-1,-2) -- (-1,2);

    \draw [red, very thick] (-3.5,1) -- (-1,1)
        -- (1,1.5) -- (3.4,0);
    \draw [help lines,dashed] (-1,1) -- (1.8,1);

    \draw (1.8,1.5) -- (1.8,0)
        node [below left] {\(P^\prime\)} -- (1.8,-1.5);
    \fill (3.4,0) circle (2pt) node [below] {\(F^\prime\)};

    \draw [|<->|] (1.8,-1.7) -- (3.4,-1.7)
        node [midway,fill] {\(f^\prime\)};

    \draw [<->] (-4,0) -- (4,0) node [right] {Eje óptico};
\end{tikzpicture}
    \end{minted}
\end{onlyenv}
\begin{onlyenv}<2->
    \begin{center}
        \begin{tikzpicture}
            \draw [cyan,<->,very thick] (1,-2) -- (1,2);
            \draw [cyan,>-<,very thick] (-1,-2) -- (-1,2);

            \draw [red, very thick] (-3.5,1) -- (-1,1) -- (1,1.5) -- (3.4,0);
            \draw [help lines,dashed] (-1,1) -- (1.8,1);

            \draw (1.8,1.5) -- (1.8,0) node [below left] {\(P^\prime\)} -- (1.8,-1.5);
            \fill (3.4,0) circle (2pt) node [below] {\(F^\prime\)};

            \draw [|<->|] (1.8,-1.7) -- (3.4,-1.7) node [midway,fill=black] {\(f^\prime\)};

            \draw [<->] (-4,0) -- (4,0) node [right] {Eje óptico};
        \end{tikzpicture}
    \end{center}
\end{onlyenv}
\end{frame}

\begin{frame}[fragile]
\frametitle{Diagrama de Cuerpo Libre}
\begin{onlyenv}<1>
    \footnotesize
    \begin{minted}{latex}
\begin{tikzpicture}
    \draw [rotate=30,thick] (-4,-0.5) -- (4,-0.5);
    \pattern [rotate=30, pattern=dots] (-4,-0.5)
        rectangle ++(8,-0.2);

    \draw [->,very thick,green] (0,0) -- (0,-2) node [below] {\(mg\)};
    \draw [rotate=30,->,very thick,orange] (0,0) -- (1,0)
        node [right] {\(F_r\)};
    \draw [rotate=30,->,very thick,red] (0,0) -- (0,1.8)
        node [above] {\(N\)};

    \filldraw [rotate=30,thick,fill=white] (-0.5,-0.5) rectangle ++(1,1);
\end{tikzpicture}
    \end{minted}
\end{onlyenv}
\begin{onlyenv}<2->
    \begin{center}
        \begin{tikzpicture}
            \draw [rotate=30,thick] (-4,-0.5) -- (4,-0.5);
            \pattern [rotate=30, pattern color=white, pattern=dots] (-4,-0.5) rectangle ++(8,-0.2);

            \draw [->,very thick,green] (0,0) -- (0,-2) node [below] {\(mg\)};
            \draw [rotate=30,->,very thick,orange] (0,0) -- (1,0) node [right] {\(F_r\)};
            \draw [rotate=30,->,very thick,red] (0,0) -- (0,1.8) node [above] {\(N\)};

            \filldraw [rotate=30,thick,fill=black] (-0.5,-0.5) rectangle ++(1,1);
        \end{tikzpicture}
    \end{center}
\end{onlyenv}
\end{frame}




\begin{frame}
\frametitle{Logo}


\begin{center}
    \lmr
    \resizebox{0.58\linewidth}{!}{
        \begin{tikzpicture}[
            star/.pic={
            \fill
                (0,-1)
                -- (0.3,-0.3)
                -- (1,0)
                -- (0.3,0.3)
                -- (0,1)
                -- (-0.3,0.3)
                -- (-1,0)
                -- (-0.3,-0.3)
                -- cycle;
            },
        ]
            \useasboundingbox (-1.5,-1.5) rectangle (1.5,1.5);

            %% Help lines
            % \begin{scope}[help lines, opacity=0.1]
            %     \draw[step=0.1] (-1.45,-1.45) grid (1.45,1.45);
            %     \draw[thin] (-1.5,0) -- (1.5,0) (0,-1.5) -- (0,1.5);
            %     \draw circle (1.5);
            % \end{scope}

            %% Planet
            \begin{scope}[rotate=30]
                %% Basic construction
                % \begin{scope}[help lines]
                %     \draw (0,0) circle [radius=0.5];
                %     \draw (0:0.025) circle [radius=0.475];
                %     \draw (0,0) circle [radius=0.45];
                %     \draw (0,0) ellipse [x radius=1, y radius=0.25];
                %     \draw (0,0) ellipse [x radius=0.95, y radius=0.20];

                %     \draw[red]
                %         (0:0.5) circle (0.5pt)
                %         arc (0:330:0.5) circle (0.5pt)
                %         to [out=60, in=5.2] (110:1 and 0.25) circle (0.5pt)
                %         arc (110:220:1 and 0.25) circle (0.5pt)
                %         to [out=-16.6, in=210] (-30:.15) circle (0.5pt);
                % \end{scope}

                \fill
                    (0:0.5)
                    arc (0:330:0.5)
                    to [out=60, in=5.2] (110:1 and 0.25)
                    arc (110:220:1 and 0.25)
                    to [out=-16.6, in=210] (-30:.15)
                    coordinate (tip)
                    to [out=207, in=-14.1] (220:0.95 and 0.20)
                    arc (220:110:0.95 and 0.20)
                    to [out=4.4, in=60] (330:0.45)
                    arc (330:180:0.45)
                    arc (180:0:0.475)
                    -- cycle;
            \end{scope}

            %% Pen nib
            \begin{scope}[shift={(tip)}, rotate=-30]
                \fill
                    (0.015,0.07)
                    -- (0.02,0.14)
                    arc (0:180:0.02)
                    -- (-0.015,0.07)
                    -- cycle;
                \fill
                    (0.003,0)
                    to[out=70, in=-30, in looseness=1.5] (0.015,0.06)
                    -- (-0.015,0.06)
                    to[out=210, out looseness=1.5, in=110]
                    (-0.003,0)
                    -- (-0.003,0.03)
                    arc (250:-70:{0.003/sin(20)})
                    -- cycle;
            \end{scope}

            %% Moons and stars
            \begin{scope}
                \fill
                    (-60:0.5) circle (0.1)
                    (135:0.5) circle (0.05)
                    (-0.8,0.2) circle (0.05)
                    (0.9,0.3) circle (0.05)
                    (0.2,0.7) circle (0.02)
                    (-0.2,0.9) circle (0.1)
                    (0.6,0.8) pic [rounded corners=0.25pt, scale=0.05] {star}
                    (-1.1,0.6) pic [rounded corners=0.25pt, scale=0.05] {star}
                    (0.8,-0.3) pic [rounded corners=0.25pt, scale=0.05] {star};
            \end{scope}

            %% Text
            \begin{scope}
                \draw (0,-0.75) node {Mundo \LaTeX{}};
            \end{scope}
        \end{tikzpicture}
    }
\end{center}
\end{frame}



\begin{frame}[allowframebreaks]
\frametitle{Referencias}
\nocite{*}
\printbibliography
\end{frame}
\end{document}
